\scanpage{45}
\mspage{44}
\mspage{78}
\margin{n°382}
\[
P_x^k:C_x^r(X)\longrightarrow P_x^k(X)
\]
soit un homomorphisme ; on a les formules :
\[
P_x^k(f)+P_x^k(g)=P_x^k(f+g),\quad
\lambda.P_x^k(f)=P_x^k(\lambda.f),\quad
P_x^k(f).P_x^k(g)=P_x^k(f.g),
\]
pour $f,g\in C_x^r(X)$, $\lambda\in K$.

7.1.3. Soit de plus h un morphisme dans une variété Z de classe $C^r$ d'un voisinage de y. Alors $j_x^k(h\circ f)$ ne dépend que de $j_x^k(f)$ et de $j_y^k(h)$ et on l'appelle le \uline{produit} (ou le \uline{composé}) de $j_x^k(f)$ et $j_y^k(h)$ ; le composé $v\circ u$ de deux jets u et v est donc défini si le but de u coïncide avec la source de v ; ce produit est associatif. Un jet u d'ordre k de X dans Y est \uline{inversible} s'il existe un jet $v\in J^k(Y,X)$ tel que $v\circ u$ et $u\circ v$ soient définis et égaux aux jets des applications identiques $1_X$ et $1_Y$. Si $k=0$, tout jet est inversible ; si $k\geq1$, un jet u est inversible si et seulement si c'est la partie principale d'ordre k d'un isomorphisme local de X dans Y ; dans ce cas, tout morphisme f tel que $j_x^k(f)=u$ est un isomorphisme local en x. Un jet d'ordre k inversible de X dans $K^p$ (resp. de source x) est appelé un \uline{repère d'ordre k} de X (resp. de source x). L'ensemble de ces repères est noté $R^k(X)$ (resp. $R_x^k(X)$). Si U et V sont des ouverts de X et Y, on a des inclusions canoniques
\[
J^k(U,V)\subset J^k(X,Y)\qquad\text{et}\qquad R^k(U)\subset R^k(X).
\]

7.1.4. Si $k'\leq k$, le jet $j_x^{k'}(f)$ d'ordre k' d'un germe de morphisme f ne dépend que de $j_x^k(f)$ et est inversible si $j_x^k(f)$ l'est. En faisant correspondre $j_x^{k'}(f)$ à $j_x^k(f)$, on définit des applications canoniques surjectives
\[
J_x^k(X,Y)_y\longrightarrow J_x^{k'}(X,Y)_y,\qquad
R_x^k(X)\longrightarrow R_x^{k'}(X)
\]
et $P_x^k(X)\longrightarrow P_x^{k'}(X)$ ; la dernière est un homomorphisme d'algèbres.

7.1.5. Soit $f=(f^1,\ldots,f^q)$ une application de classe $C^r$ d'un ouvert U de $K^p$ dans l'espace $K^q$. Si $x\in U$, le jet $j_x^k(f)$ de f est complètement déterminé par les coordonnées de sa source x et par les valeurs en x des fonctions $\Delta^\alpha f^i$, avec
